% =============================================================================
% INTRODUCCIÓN
% =============================================================================

\chapter*{INTRODUCCIÓN}
\addcontentsline{toc}{chapter}{INTRODUCCIÓN}
\markboth{INTRODUCCIÓN}{}

La evaluación por competencias se ha consolidado como un eje estructurante de la transformación educativa contemporánea, desplazando el énfasis desde la mera memorización de contenidos hacia la capacidad de movilizar conocimientos, habilidades y actitudes en situaciones auténticas \cite{ref01_evalcompetencias,ref33_cano}. En el contexto de la educación básica elemental, este enfoque adquiere una relevancia particular, dado que los primeros años de escolaridad constituyen el cimiento sobre el cual se erige el desarrollo cognitivo posterior, especialmente en disciplinas instrumentales como las matemáticas \cite{ref35_pisa,ref36_niss}.

A pesar de los avances teóricos, la realidad de las aulas ecuatorianas evidencia una brecha significativa entre los postulados pedagógicos y los instrumentos disponibles para operacionalizarlos. Las herramientas tecnológicas existentes para la evaluación por competencias, como SCALA \cite{ref21_scala_universidad}, CAT \cite{ref10_cat}, SECAT \cite{ref11_secat} y MCAT, fueron concebidas mayoritariamente para la educación superior, dejando desatendido el nivel básico elemental. Este vacío instrumental obliga a los docentes de primaria a recurrir a métodos tradicionales que no logran capturar de forma fiel el desarrollo de las destrezas matemáticas críticas, ni atender la diversidad cognitiva propia de la infancia.

Frente a esta problemática, el presente proyecto tecnológico se propuso desarrollar \textbf{Kiddy Quiz}, una aplicación web inclusiva orientada al control y evaluación de competencias curriculares en matemáticas para estudiantes de educación básica elemental. El diseño de la solución se sustentó en una investigación de campo que incluyó entrevistas semiestructuradas con cinco docentes con experiencia en evaluación por competencias y dos psicopedagogos, lo cual permitió identificar las barreras cognitivas y de ansiedad que afectan a los niños durante los procesos evaluativos tradicionales.

El desarrollo del software se gestionó mediante la metodología ágil de Programación Extrema (XP) \cite{ref53_beck,ref47_ganney_xp}, adaptada al contexto de un desarrollo individual mediante ciclos planificados de retroalimentación con expertos. La arquitectura del sistema se construyó utilizando Angular \cite{ref49_angular} para el \textit{frontend}, NestJS sobre Node.js \cite{ref50_nodejs} para el \textit{backend} y PostgreSQL \cite{ref51_postgresql} como motor de base de datos relacional. Como innovaciones centrales del proyecto destacan: la integración de interfaces basadas en pictogramas para favorecer la navegación autónoma infantil y la implementación del modelo de inteligencia artificial generativa Gemini para producir retroalimentación motivacional adaptativa tras cada evaluación.

La validación del producto se realizó mediante la Escala de Usabilidad del Sistema (SUS) y el Modelo de Aceptación Tecnológica (TAM), aplicados a quince estudiantes y cinco docentes. Los resultados evidenciaron una usabilidad en categoría \textit{Excelente} y una intención de uso prácticamente unánime entre los estudiantes, validando empíricamente la pertinencia pedagógica y técnica de la solución.

El documento se estructura en cinco capítulos. El \textbf{Capítulo I} presenta el marco contextual del proyecto, incluyendo el planteamiento del problema, los objetivos y la justificación. El \textbf{Capítulo II} desarrolla el marco teórico, articulando un marco conceptual sobre evaluación por competencias, tecnología educativa y metodologías ágiles, junto con un marco referencial que analiza herramientas existentes en el dominio. El \textbf{Capítulo III} expone la metodología, describiendo las fases de desarrollo, el cronograma, los recursos y los indicadores de evaluación. El \textbf{Capítulo IV} presenta los resultados, abarcando el producto tecnológico desarrollado y la comprobación empírica mediante los instrumentos SUS y TAM. Finalmente, el \textbf{Capítulo V} sintetiza las conclusiones del estudio y formula recomendaciones orientadas al fortalecimiento futuro de la solución y a su escalamiento institucional.
